\section{Đánh giá Hệ thống}

Chương này trình bày quá trình đánh giá toàn diện hệ thống theo từng thành phần trong pipeline, từ xử lý tín hiệu tiếng nói đến truy xuất và tổng hợp thông tin. Mục tiêu của việc đánh giá là kiểm chứng tính chính xác, độ ổn định và khả năng hoạt động trong điều kiện dữ liệu thực tế có nhiễu.

Do hệ thống được xây dựng theo kiến trúc nhiều tầng, quá trình đánh giá được tổ chức theo hướng phân rã theo từng mô-đun, kết hợp với các chỉ số định lượng và các kịch bản thực nghiệm.